% Einheit 2 von 3 – Mit der Näherungsformel rechnen: die Grenzgleichung nach p lösen (quadratisch (bank/konfidenzintervalle/e2.jsonl, zusammenbau v0.1)
%% TODO Zeitmarke Einheit 2: Marken-Zeile nicht lesbar
%% TODO Zweigzeile Teil 1: was hier gelernt wird (Satz in Schülersprache aus dem Einheitstitel) – Einheit 2
\einheitenkopf[e2]{Einheit 2 von 3 · Mit der Näherungsformel rechnen: die Grenzgleichung nach p lösen (quadratisch}
\zweigzeile{keine Prüfungsaufgabe}
\setcounter{aufgabe}{10}

%% TODO Titel: Ich-kann-Satz fehlt in der Bank (Platzhalter: Kettenname) – Nr. 11
%% TODO Anweisung: ein Satz über der ersten Teilaufgabe fehlt in der Bank – Nr. 11
\begin{aufgabe}{Näherungsformel}
%% TODO \rechenplatz{n}: Schrittzahl des Grundfalls fehlt in der Bank (nur bei fester Schreibform, 2.3 b)
\begin{teile}
\teil Vorwärts oder rückwärts? Nur ankreuzen: „Von 500 Befragten sagen 140 ja. Bestimme das Konfidenzintervall.“ \\ \kreuz{vorwärts: Grenzen gesucht} \\ \kreuz{rückwärts: Ergebnis oder Umfang gesucht}
\end{teile}
\begin{gleichungsraster}[2]
\gl{\text{In einer Stichprobe vom Umfang 400 gibt es 120 Treffer. Berechne mit der Grenzgleichung die Grenzen des Konfidenzintervalls zur Sicherheitswahrscheinlichkeit 95\,\% auf Tausendstel.}} &
\gl{\text{Die untere Grenze eines Konfidenzintervalls zur Sicherheitswahrscheinlichkeit 95\,\% ist $0{,}319$, der Umfang 500. Bestimme die Trefferzahl der Stichprobe.}} \\
\gl{\text{Bei einem Glücksrad nimmt man an, dass das Feld „Stern“ mit der Wahrscheinlichkeit $0{,}2$ kommt. Beobachtet wird, dass genau 25\,\% der Drehungen „Stern“ ergeben. Ermittle die kleinste Anzahl von Drehungen, bei der dieser Anteil zur Sicherheitswahrscheinlichkeit 95\,\% nicht mit $0{,}2$ verträglich ist (Abitur 2023 LK).}} & \\
\end{gleichungsraster}
\end{aufgabe}

%% TODO Titel: Ich-kann-Satz fehlt in der Bank (Platzhalter: Kettenname) – Nr. 12
%% TODO Anweisung: ein Satz über der ersten Teilaufgabe fehlt in der Bank – Nr. 12
\begin{aufgabe}{Verträglichkeit einer vermuteten Trefferwahrscheinlichkeit über das 1,96σ-Intervall um den Erwartungswert prüfen}
\begin{teile}
\teil Für Blumenzwiebeln wird eine Austriebswahrscheinlichkeit von 70\,\% vermutet; von 80 gesetzten Zwiebeln treiben 50 aus. Die Vermutung gilt zur Sicherheitswahrscheinlichkeit 95\,\% als verträglich, wenn 50 im Intervall von $\mu - 1{,}96\sigma$ bis $\mu + 1{,}96\sigma$ liegt. Untersuche das \hfill (Abitur 2017 LK).
\end{teile}
\end{aufgabe}

%% TODO Titel: Ich-kann-Satz fehlt in der Bank (Platzhalter: Kettenname) – Nr. 13
%% TODO Anweisung: ein Satz über der ersten Teilaufgabe fehlt in der Bank – Nr. 13
\begin{aufgabe}{Näherungsformel – Fehler finden, Begründen}
\begin{teile}
\teil Zu 90 Treffern bei Umfang 300 soll das Konfidenzintervall zu 95\,\% bestimmt werden. Jan rechnet $0{,}3 \pm 1{,}96 \cdot \sqrt{0{,}3 \cdot 0{,}7 / 300}$ und erhält $[0{,}248; 0{,}352]$. Finde den Fehler und rechne richtig.
\teil Begründe, warum die Grenzgleichung zwei Lösungen für p hat und warum man beide braucht.
\end{teile}
\end{aufgabe}
